\section{Future Work}\label{sec:future}

\textbf{Additional architectures.}
\toolname currently targets x86-64 binaries.
\QEMU supports many additional architectures (ARM, MIPS, RISC-V, PowerPC), and the only architecture-dependent component is the trace-log parser.
Extending to other architectures requires adapting the regular expression that extracts basic-block boundaries from \QEMU's log output.

\textbf{Differential fuzzing guided by labels.}
Once \toolname has labeled a parser's CFG, the labels provide a natural coverage metric for format-aware fuzzing.
A fuzzer could prioritize inputs that exercise under-explored grammar productions, using \toolname's labels as a coverage map.

\textbf{Grammar mining integration.}
As discussed in Section~\ref{sec:rel:grammar}, grammar-recovery tools and \toolname solve complementary problems.
An integrated pipeline could first recover an approximate grammar from a binary, then use \toolname to validate and refine the mapping, closing the loop between ``what does this binary eat?'' and ``where does it eat it?''

\textbf{Hierarchical labeling.}
The current template system supports one level of production-to-label mapping.
Extending to hierarchical templates (labeling not just ``string'' but ``string-escape-sequence'' versus ``string-literal-characters'') would increase the granularity of the output without algorithmic changes.

\section{Ethical Considerations}\label{sec:ethics}

\toolname is a defensive reverse-engineering tool that annotates binary code with format-semantic labels.
It does not discover, create, or exploit vulnerabilities.
All evaluated targets are publicly available software projects; the artifact contains no proprietary or classified binaries.
The present repository does not include a license file, so it should be treated as source-available for research review rather than as a licensed open-source release.

\section{LLM Usage Statement}\label{sec:llm}

Large language models were used for editorial assistance and exploratory checks while preparing this draft.
The committed ELF label counts reported here were manually reviewed; no unverified model-generated measurement is presented as a final result.

\section{Conclusion}\label{sec:conclusion}

We presented \toolname, a technique for associating known input-format labels with locations in an unknown binary.
By systematically generating structurally varied inputs, recording their \QEMU translation logs, and associating record-pair deltas with parse-tree deltas, \toolname produces candidate semantic annotations without requiring source code, compiler instrumentation, or data-flow analysis.

The artifact contains completed runs for URI, JSON, and ELF; manual validation is complete only for ELF, where 62 of 68 labels are strict true positives and one additional label is marginal.
These results support continued study, but broader accuracy claims require review of the URI and JSON outputs, completion of the planned APR and HPACK runs, and replacement or validation of the translation-log tracer with dynamic execution events.
